\documentclass[../DoAn.tex]{subfiles}
\begin{document}

\begin{center}
    \Large{\textbf{TÓM TẮT NỘI DUNG ĐỒ ÁN}}\\
\end{center}
\vspace{1cm}
Trong bối cảnh chuyển đổi số giáo dục, học sinh và giáo viên Trung học Phổ thông (đặc biệt môn Tin học) đối mặt với nhiều thách thức. Học sinh cần môi trường học tương tác, còn giáo viên mất nhiều thời gian soạn giáo án, sinh bài tập và chấm điểm. Mặc dù các mô hình ngôn ngữ lớn (LLM) mang lại tiềm năng lớn, chúng vẫn gặp rào cản về việc sinh thông tin sai lệch (hallucination) và không bám sát chương trình chuẩn. Từ đó, bài toán đặt ra là xây dựng một trợ lý ảo giáo dục thông minh có khả năng xử lý tự động, chính xác các tác vụ dựa trên nguồn tri thức chuẩn mực.

Nghiên cứu ứng dụng kiến trúc tìm kiếm và sinh văn bản (RAG) kết hợp phương pháp điều phối đa tác nhân (Multi-Agent LLM). Cụ thể, cơ sở dữ liệu vector được xây dựng từ sách giáo khoa chuẩn. Quá trình truy xuất dữ liệu áp dụng tìm kiếm lai (từ khóa và ngữ nghĩa), tối ưu bằng thuật toán Reciprocal Rank Fusion và mạng Cross-Encoder để xếp hạng lại nhằm trích xuất chính xác ngữ cảnh. Đồng thời, hệ thống bóc tách, đánh giá và định tuyến đa ý định thông qua bộ điều khiển luồng trung tâm (Pipeline Controller). Các dịch vụ miền sau đó xử lý lệnh song song cho các tác vụ như hỏi đáp, sinh bài tập đa định dạng, chấm điểm tự luận theo tiêu chí và tạo cấu trúc bài giảng.

Thử nghiệm trên 246 mẫu kiểm thử bằng khung đo lường RAGAS cho kết quả rất khả quan. Hệ thống đạt tính trung thực (Faithfulness) 0.975 giúp khắc phục gần như hoàn toàn hiện tượng ảo giác, độ phủ ngữ cảnh (Context Recall) đạt 0.959 và mức độ liên quan (Answer Relevancy) đạt 0.857. Tốc độ phản hồi trung bình của toàn hệ thống duy trì ổn định ở mức 6,5 giây cho mỗi truy vấn phức tạp.

Nghiên cứu cung cấp một giải pháp công nghệ giáo dục (EdTech) hiệu quả, hỗ trợ tự động hóa một số quy trình dạy và học môn Tin học THPT. Đo thời, kết quả chứng minh sự thành công của kiến trúc lai RAG và Multi-Agent trong việc định hướng LLM, tạo nền tảng vững chắc để phát triển các hệ sinh thái học tập cá nhân hóa sâu rộng hơn.

\begin{flushright}
Sinh viên thực hiện\\
\begin{tabular}{@{}c@{}}
\textit{(Ký và ghi rõ họ tên)}
\end{tabular}
\end{flushright}

\end{document}